\documentclass[11pt,a4paper]{article}
\usepackage[margin=1in]{geometry}
\usepackage{booktabs}
\usepackage{amsmath,amssymb}
\usepackage{graphicx}
\usepackage{hyperref}
\usepackage{xcolor}
\usepackage{enumitem}
\usepackage{caption}
\usepackage{subcaption}
\usepackage{float}
\usepackage{multirow}
\usepackage[T1]{fontenc}
\usepackage{times}
\usepackage{microtype}

\hypersetup{colorlinks=true, linkcolor=blue!60!black, citecolor=blue!60!black, urlcolor=blue!60!black}

\title{\textbf{Baddie: Skeleton-Based Action Recognition in Badminton}}
\author{}
\date{}

\begin{document}
\maketitle

\begin{abstract}
Fine-grained shot-type classification in badminton is difficult because many shot types share near-identical body poses and differ only in racket speed or shuttle trajectory---information that 2D skeleton representations capture poorly. We present an end-to-end pipeline that extracts dual-player skeletons from broadcast video using Grounding DINO-guided YOLOv8-Pose, constructs court-normalised spatio-temporal graphs with hierarchical feature engineering, and classifies 17 shot types via ST-GCN. Through sequential ablations we show that (1) shuttlecock trajectory fusion via cross-attention is the single most impactful design choice (+4.7 macro-F1), (2) the opponent's skeleton provides meaningful court-positioning context (+4 F1 over single-player), and (3) the dominant performance bottleneck is biomechanical ambiguity between visually similar shot types, not model capacity. Our best configuration achieves 0.593 macro-F1 and 0.906 top-3 accuracy on held-out test matches, with leave-2-out cross-validation confirming $\pm$0.09 F1 match-level variance across folds.
\end{abstract}

% ============================================================================
\section{Introduction}

Tactical analysis in badminton remains a manually intensive process. Coaches review match footage frame-by-frame to identify shot types, assess execution quality, and study opponent patterns. This process is time-prohibitive, subjective, and inaccessible to amateur and semi-professional players who lack dedicated analyst support.

Automated shot-type classification from video is challenging because it requires understanding both player body mechanics and shuttlecock trajectory---information that is spatially distributed across the court and temporally spread across the shot sequence. Existing approaches fall into three broad categories. Frame-level visual methods~\cite{shuttlenet} classify shots directly from video frames using convolutional features, achieving reasonable accuracy but requiring high computational cost and exhibiting sensitivity to camera angle, lighting, and jersey colour. Sequence-level metadata methods~\cite{shuttleset} use rally-level features (shot order, player position coordinates, shuttle landing zones) to forecast future shots, bypassing the visual recognition problem entirely but requiring ground-truth annotations at inference time. Skeleton-based methods~\cite{tempose, bst} extract player joint coordinates and classify from pose sequences, offering a middle ground: compact, camera-invariant representations that capture biomechanics while remaining tractable for graph and transformer architectures.

The most relevant prior work is BST (Badminton Stroke-type Transformer)~\cite{bst}, which combines skeleton joints, bone vectors, shuttlecock trajectory, and player court positions through a transformer architecture with specialised trajectory-cleaning modules. BST achieves 0.770 accuracy and 0.704 macro-F1 on ShuttleSet using 40 matches and 35 classes (17 shot types $\times$ 2 court sides). TemPose~\cite{tempose} uses a similar transformer backbone with early fusion of skeleton, shuttle, and position modalities.

Our work differs in three respects. First, we address the skeleton extraction quality problem directly: standard YOLOv8-Pose consistently detects non-player entities (chair umpires, audience members) with high confidence, contaminating downstream representations. We introduce a Grounding DINO-guided extraction pipeline with sport-specific text prompting that substantially reduces this contamination. Second, we use a graph convolutional architecture (ST-GCN) rather than a transformer, exploiting the known anatomical connectivity of the skeleton as an inductive bias rather than learning attention patterns from data. Third, we conduct systematic sequential ablations across skeleton representation, temporal windowing, shuttle fusion, and classifier head to isolate the contribution of each design choice---revealing that shuttle trajectory fusion via cross-attention contributes more than any architectural modification.

% ============================================================================
\section{Datasets and Data Acquisition}

\subsection{ShuttleSet}

The ShuttleSet dataset~\cite{shuttleset} contains annotated CSV files for 44 professional broadcast matches with 17 shot-type categories, player positions, and shuttlecock coordinates. From the original 44 matches, 6 were excluded due to broken YouTube links or non-standard camera angles (side court or overhead). Women's matches and matches featuring left-handed players were excluded to reduce confounding variation from fundamentally different pose distributions at the current data scale.

The remaining 21 matches (approximately 18,000 labelled shots) were split into 17 training matches ($\sim$12,300 labelled shots after skeleton extraction filtering), 2 validation matches, and 2 held-out test matches. Frames were downloaded at every 4th frame plus every annotated hit frame, yielding approximately 10 FPS---balancing storage efficiency with sufficient temporal resolution to track player movements.

\paragraph{Class imbalance.} The most frequent shot type (net drop, 18.4\%) has approximately 29$\times$ more training samples than the rarest well-represented class. The top-5 types (net drop, lob/lift, block, push, clear) account for approximately 60\% of all shots. We address this through class-weighted cross-entropy loss during training, though the imbalance is not fully compensated as shown in the per-class analysis (Section~\ref{sec:perclass}).

% ============================================================================
\section{Dataset Preprocessing}

From extracted frames we derive three signal types: player skeletons, court-normalised positions via per-match homography, and shuttlecock trajectories. A manual annotation QA tool was developed to verify extraction accuracy, allowing visual inspection of skeleton overlays, player movement, and shot-type labels across rallies.

\subsection{Skeleton Extraction}

Skeleton keypoints were extracted using YOLOv8-Pose (yolov8s-pose), which performs joint person detection and 17-joint COCO keypoint estimation in a single forward pass. A systematic quality audit revealed significant contamination: YOLOv8-Pose imposes no spatial constraint on detections, so the chair umpire---stationary, upright, and well-lit---is consistently detected with high confidence. Audience members were also occasionally detected.

Two corrections were implemented:
\begin{enumerate}[leftmargin=*]
    \item \textbf{Grounding DINO-guided extraction:} Grounding DINO (grounding-dino-tiny) with the text prompt ``badminton player'' produces court-region bounding boxes. Only YOLOv8 detections with IoU $\geq 0.25$ against a GDINO box are accepted. If fewer than 2 valid players pass this filter, the pipeline falls back to plain YOLOv8 top-2 selection.
    \item \textbf{Pixel boundary constraint:} Since broadcast cameras keep the court centred, valid detections are restricted to $X_{\min} = 450$, $X_{\max} = 1450$ pixels, excluding detections outside the court region.
\end{enumerate}

\subsection{Shuttlecock Extraction}

TrackNetV3 was used to extract shuttlecock trajectories on a per-rally basis, providing $(x, y)$ coordinates and detection confidence at each frame. Velocity channels $(dx, dy)$ are computed via finite differences during feature construction.

\subsection{Player Court Ordering and Hitter Identification}

Shot-type identification from skeletal poses requires knowledge of which player is the hitter and their court position, because a player facing toward versus away from the camera produces fundamentally different skeleton appearances---joint ordering, occlusion patterns, and limb foreshortening are all affected by the perspective.

After selecting the top-2 detections per frame, both players are sorted by their mean Y-coordinate centroid. Since image Y increases downward, the player with smaller mean Y is the top-court player ($P_0$, joints 0--16) and the player with larger mean Y is the bottom-court player ($P_1$, joints 17--33). This assignment is computed purely from image geometry.

The hitter is then identified via the annotated \texttt{player\_location\_y} field from ShuttleSet metadata. Per-rally median Y values for players A and B determine the top/bottom mapping, computed per-rally because players switch sides between games. Each training sample carries a \texttt{hitter} field of either ``top'' or ``bottom''.

Note that BST~\cite{bst} handles the court-side distinction by doubling the label space to 35 classes (17 types $\times$ 2 sides + ``none''). We instead retain 17 classes and provide court-side information through the consistent player ordering in the skeleton graph ($P_0$ = top, $P_1$ = bottom) and an optional hitter identity channel (tested in ablation A5--A6). This avoids halving per-class sample counts, which is particularly important at our data scale.

\subsection{Shot Window Selection}

Each training sample is a temporal window of skeleton sequences centred on the annotated hit frame. Frames before the hit capture the preparation phase (backswing, footwork), where shot types become most distinguishable. Frames after capture the follow-through and recovery.

We selected $T=32$ frames ($\pm$16 around the hit frame) based on the distribution of inter-shot intervals, covering approximately the 75th percentile. $T=16$ is too short for full backswing and follow-through; $T=50$ risks significant overlap with adjacent shots. Overlapping windows between consecutive shots within a rally are expected---each sample has a different hit frame, label, and potentially different hitter, so overlap does not inflate evaluation metrics.

\subsection{Feature Engineering}

Raw keypoint coordinates are not sufficiently expressive to distinguish shot types: a player's wrist may occupy the same position for both a drop shot and a smash---what differs is velocity, court position, and arm angle. We derive richer signals through cumulative feature layers:

\begin{table}[H]
\centering
\caption{Feature layer hierarchy. Each layer includes all features from previous levels.}
\label{tab:features}
\begin{tabular}{@{}llp{8cm}@{}}
\toprule
\textbf{Layer} & \textbf{Dims} & \textbf{Description} \\
\midrule
L0: Position & 2 & Raw pixel coordinates $[x, y]$ \\
L1: Kinematics & +4 & Velocity and acceleration $[v_x, v_y, a_x, a_y]$ via finite differences \\
L2: Court context & +3 & Distances to net, court centre, and opponent centroid (metres, via homography) \\
Hitter flag & +1 & 1.0 for hitter's joints, 0.0 for opponent's joints \\
L3: Biomechanical & +3 & Joint angles at elbow, shoulder, and knee \\
\bottomrule
\end{tabular}
\end{table}

Homography is applied only for L2 court-context features, where physical distances in metres are needed. The camera's elevated viewing angle introduces perspective distortion: a player near the camera (bottom court) covers more pixels per metre of movement than a player far from the camera. Per-match homography matrices were computed from court boundary annotations (via Roboflow) to transform player centroids from pixel to court-metre coordinates. L0, L1, and L3 features remain in pixel space to preserve body proportions and motion patterns.

% ============================================================================
\section{Encoder: ST-GCN}

The enriched feature tensor is processed by a Spatial-Temporal Graph Convolutional Network (ST-GCN)~\cite{stgcn}. Each of the 9 blocks performs two operations:
\begin{itemize}[leftmargin=*]
    \item \textbf{Spatial graph convolution:} Aggregates information along the skeleton topology (e.g., wrist $\to$ elbow $\to$ shoulder), learning biomechanically meaningful patterns within each frame. The graph structure is fixed by anatomical connectivity, providing an inductive bias well-suited for skeleton data where joint relationships are known \textit{a priori}.
    \item \textbf{Temporal convolution:} A 1D convolution with kernel size 9 slides across the time axis, capturing how spatial relationships evolve---for example, the arm extending rapidly over several frames indicates a smash, while slow controlled motion suggests a drop shot.
\end{itemize}

The architecture uses progressive channel expansion (64 $\to$ 128 $\to$ 256) with strided temporal convolutions at blocks 3 and 6 that compress the temporal dimension ($32 \to 16 \to 8$ frames), progressively building from local joint movements to full-body coordination patterns. Global average pooling at the final block produces a 256-dimensional embedding. Dropout of 0.3 is applied throughout.

\paragraph{Architecture choices.} The 9-layer depth, 256-dimensional embedding, and kernel size 9 follow the original ST-GCN~\cite{stgcn} defaults for 17-joint COCO skeletons. We did not tune these hyperparameters, prioritising systematic ablation of input representation and fusion strategies over architecture search. The temporal kernel of 9 frames corresponds to approximately 0.9 seconds at our 10 FPS extraction rate, covering the core phase of most shot executions.

\paragraph{Comparison to transformer approaches.} BST~\cite{bst} and TemPose~\cite{tempose} use transformer encoders that apply self-attention over the temporal dimension. Transformers can capture long-range dependencies without the locality constraint of convolutional kernels, but require learning spatial relationships from data rather than encoding them via graph structure. With $\sim$12,000 training shots, the graph inductive bias of ST-GCN provides useful structure that a transformer would need to learn from limited data. A Skeleton Transformer encoder was implemented but not evaluated due to scope constraints; direct comparison remains future work.

% ============================================================================
\section{Supervised Training: Sequential Ablation Study}

We perform sequential ablations on 17-class shot-type classification to isolate the contribution of each design choice. Each group builds on the winner of the previous group. All experiments use AdamW (lr = $10^{-3}$), cosine annealing, class-weighted cross-entropy, and early stopping (patience = 10) on validation macro-F1. Training uses 17 matches ($\sim$12,300 labelled shots), with 2 matches for validation and 2 for held-out test. We report macro-F1, accuracy, and top-$k$ accuracy on the test set (1,627 labelled shots).

The sequential design reflects a dependency structure: skeleton representation (Group~A) determines what the model sees; temporal windowing (Group~B) determines how much context it processes; shuttle fusion (Group~C) adds an auxiliary modality that interacts with both; the classifier head (Group~D) operates on the encoder output and cannot compensate for upstream deficiencies.

\subsection{Data Augmentation}

All training runs apply skeleton-level augmentations on-the-fly: (i)~joint coordinate jittering ($\sigma = 0.02$); (ii)~temporal crop and resample ($\pm$20\% window variation, resampled to $T=32$); (iii)~spatial rotation ($\pm$15$^\circ$ around court centre); and (iv)~joint masking (10--20\% of joints zeroed per frame). These augmentations address the severe overfitting observed (train loss reaches near-zero while validation F1 plateaus at $\sim$0.55), preventing memorisation of specific player proportions, shot timings, and individual joint trajectories.

\subsection{Group A: Skeleton Representation}

We vary three axes: feature layer (L2 vs L3), player configuration (dual 34-node graph vs single 17-node hitter-only), and hitter identity (binary channel indicating the striking player).

\begin{table}[H]
\centering
\caption{Group A results. Fixed $T=32$, no shuttle, linear head.}
\label{tab:groupA}
\begin{tabular}{@{}lcccc@{}}
\toprule
\textbf{Ablation} & \textbf{Acc.} & \textbf{F1} & \textbf{Top-3} & \textbf{Top-5} \\
\midrule
A1: Dual-player, L2 (9-dim) & 0.606 & \textbf{0.546} & 0.861 & 0.912 \\
A2: Dual-player, L3 (12-dim) & 0.553 & 0.522 & 0.896 & 0.950 \\
A3: Single-player, L2 (9-dim) & 0.576 & 0.504 & 0.868 & 0.938 \\
A4: Single-player, L3 (12-dim) & 0.571 & 0.516 & 0.865 & 0.937 \\
A5: Dual-player, L2 + hitter (10-dim) & 0.626 & 0.544 & 0.882 & 0.946 \\
A6: Dual-player, L3 + hitter (13-dim) & 0.566 & 0.498 & 0.883 & 0.943 \\
\bottomrule
\end{tabular}
\end{table}

\textbf{Winner: A1 (dual-player, L2).} Three findings emerge:

\textit{Dual outperforms single} by approximately 4 macro-F1 points (A1 vs A3). The opponent's skeleton provides court-positioning context: if the opponent is at the net, the hitter is more likely executing a lob or clear; if the opponent is at the baseline, a drop or net shot becomes more probable. The inter-player graph edges in ST-GCN carry meaningful tactical signal, consistent with badminton's relational spatial dynamics.

\textit{L2 outperforms L3} consistently. The additional joint-angle features in L3 add 3 dimensions per joint without improving generalisation. L3 models early-stop at 26--28 epochs versus 72 for A1, indicating accelerated overfitting. The court-context distances in L2 appear to be the most informative features beyond raw position and kinematics.

\textit{The hitter channel is inconclusive.} A5 achieves the highest accuracy (0.626) but marginally lower F1 than A1 (0.544 vs 0.546). The hitter label may help on easy classes (serve identification, where court side is informative) without improving the hard classes that dominate macro-F1. A6 overfits substantially (validation--test gap of +0.08 F1).

\subsection{Group B: Temporal Window}

\begin{table}[H]
\centering
\caption{Group B results. Using Group A winner (dual-player, L2).}
\label{tab:groupB}
\begin{tabular}{@{}lcc@{}}
\toprule
\textbf{Ablation} & \textbf{Acc.} & \textbf{F1} \\
\midrule
B1: Fixed $T=32$ & \textbf{0.606} & \textbf{0.546} \\
B2: Variable (hit-to-hit, resampled) & 0.558 & 0.476 \\
\bottomrule
\end{tabular}
\end{table}

\textbf{Winner: B1 (fixed window).} The variable window degrades by 7 F1 points. Short inter-hit intervals (e.g., 10 frames during fast net exchanges) are stretched by 3$\times$, manufacturing synthetic motion; long intervals (60+ frames) are compressed, discarding motion phases. The ST-GCN's fixed temporal kernel (size 9) expects consistent cadence---variable resampling destroys this. Frame-rate normalisation or adaptive temporal convolutions may recover the intended benefit.

\subsection{Group C: Shuttle Fusion}

Using the Group A--B winner, we test two shuttle integration strategies. The shuttle trajectory provides $(x, y, dx, dy)$ across the temporal window.

\begin{table}[H]
\centering
\caption{Group C results.}
\label{tab:groupC}
\begin{tabular}{@{}lcccc@{}}
\toprule
\textbf{Ablation} & \textbf{Acc.} & \textbf{F1} & \textbf{Top-3} & \textbf{Top-5} \\
\midrule
C1: No shuttle (baseline) & 0.606 & 0.546 & 0.861 & 0.912 \\
C2: Shuttle as graph node & 0.596 & 0.529 & 0.881 & 0.949 \\
C3: Shuttle cross-attention & \textbf{0.619} & \textbf{0.593} & \textbf{0.906} & \textbf{0.961} \\
\bottomrule
\end{tabular}
\end{table}

\textbf{Winner: C3 (cross-attention)}, the largest single improvement (+4.7 F1). A 1D temporal CNN encodes the shuttle into a 128-dim sequence, which attends to the 256-dim skeleton embedding via 4-head cross-attention. The graph-node approach (C2) hurts ($-$1.7 F1): forcing the shuttle through skeleton message-passing corrupts both representations, since the shuttle evolves ballistically while joints are biomechanically constrained. Cross-attention allows selective attendance---strong when the trajectory disambiguates (steep angle $\Rightarrow$ smash vs drop), weak when shuttle data is noisy or missing. Top-3 accuracy reaches 0.906, meaning the correct shot type is in the model's top-3 for over 90\% of samples. This aligns with BST~\cite{bst}, which similarly finds shuttle trajectory to be the most impactful auxiliary modality.

\subsection{Group D: Classifier Head}

With the encoder frozen at C3, we compare classifier heads.

\begin{table}[H]
\centering
\caption{Group D results. Encoder frozen from C3 checkpoint.}
\label{tab:groupD}
\begin{tabular}{@{}lcc@{}}
\toprule
\textbf{Ablation} & \textbf{Acc.} & \textbf{F1} \\
\midrule
D1: Linear ($256 \to 17$) & 0.619 & 0.593 \\
D2: MLP-1 ($256 \to 128 \to 17$) & \textbf{0.626} & \textbf{0.576} \\
D3: MLP-2 ($256 \to 128 \to 64 \to 17$) & 0.628 & 0.572 \\
\bottomrule
\end{tabular}
\end{table}

\textbf{Winner: D2} by validation F1 (0.594 vs 0.587 for D3). Both MLPs improve accuracy over the linear head, confirming the embedding is not fully linearly separable. Diminishing returns from D2 to D3 suggest the needed nonlinearity is modest---a single hidden layer suffices.

% ============================================================================
\section{Cross-Validation}

Leave-2-matches-out cross-validation on A1 across 10 folds:

\begin{center}
\textbf{Mean macro-F1: $0.580 \pm 0.089$} \quad (range: 0.392--0.670)
\end{center}

The high variance reflects match-level distributional shift: different player--tournament combinations exhibit distinct shot distributions. The worst fold (0.392) involves players whose patterns are poorly represented in training. Ablation differences below $\sim$0.02 F1 on a single split should be interpreted cautiously. The full cascade was not cross-validated due to computational constraints.

% ============================================================================
\section{Per-Class Diagnostic}
\label{sec:perclass}

\begin{table}[H]
\centering
\caption{Per-class F1 and primary confusion patterns (C3 model).}
\label{tab:perclass}
\begin{tabular}{@{}lcp{7.5cm}@{}}
\toprule
\textbf{Shot Type} & \textbf{F1} & \textbf{Primary Confusion} \\
\midrule
short\_serve & 0.97 & Near-perfect; distinctive pose \\
clear & 0.90 & Distinctive overhead motion \\
net\_drop & 0.81 & Benefits from large sample size ($n = 2{,}753$) \\
block & 0.66 & Some confusion with drive \\
transition & 0.54 & Confused with slice\_drop \\
lob\_lift & 0.52 & High recall (0.88) but low precision; absorbs push \\
smash & 0.49 & Bidirectional confusion with tap\_smash (36--37\%) \\
tap\_smash & 0.40 & Bidirectional confusion with smash \\
push & 0.29 & 72\% misclassified as lob\_lift \\
drive & 0.29 & Confused with push, block, tap\_smash \\
\bottomrule
\end{tabular}
\end{table}

The dominant failure mode is biomechanical ambiguity: shot types sharing near-identical poses but differing in racket speed or shuttle trajectory. Push/lob\_lift is the most extreme---both involve a forward body position with upward arm motion, distinguished by wrist angle and shuttle launch speed that 2D coordinates capture poorly. Smash/tap\_smash exhibits symmetric confusion (36--37\%), reflecting a power distinction rather than form. Shuttle cross-attention partially addresses this---push F1 improves from 0.29 to 0.49 with trajectory---but residual confusion represents a fundamental limitation of skeleton-based features for fine-grained discrimination. More discriminative shuttle features (speed at contact, launch angle, post-contact arc) may require longer temporal context or explicit velocity modelling.

% ============================================================================
\section{Summary}

\textbf{Best configuration:} Dual-player L2 + fixed $T=32$ + shuttle cross-attention + MLP-1 head: macro-F1 = 0.593, accuracy = 0.619, top-3 = 0.906.

\textbf{Shuttle trajectory is the most impactful choice} (+4.7 F1). Cross-attention outperforms graph-node fusion, consistent with the principle that modalities with different spatial semantics should be fused at the representation level.

\textbf{The bottleneck is inter-class confusion, not architecture.} Push/lob\_lift, smash/tap\_smash, and drive/block account for the majority of errors---a representation problem requiring features beyond 2D skeleton coordinates.

\textbf{Class imbalance is severe} (29$\times$). Weighted loss helps but does not fully compensate.

\textbf{Match-level variance is high} ($\pm$0.09 F1), confirming single-split uncertainty.

% ============================================================================
\section{Limitations and Future Work}

\textbf{Data scale and feature representation.} We use 21 matches versus BST's 40, and our feature engineering lacks bone vectors (2D limb direction/length) and bounding-box-normalised coordinates, both of which BST employs. Adding bone features as a parallel input stream and normalising skeleton coordinates relative to the player's bounding box are the most directly actionable improvements.

\textbf{2D skeleton limitations.} Depth information is lost, affecting shots with vertical shuttle trajectories. Monocular 3D pose lifting (MotionBERT) could recover this dimension.

\textbf{Fixed graph topology.} ST-GCN's predefined adjacency cannot learn non-anatomical relationships. Adaptive-topology variants (CTR-GCN, InfoGCN) may help but were not explored.

\textbf{Encoder comparison.} An implemented Skeleton Transformer was not evaluated. Direct comparison against ST-GCN on identical data would clarify whether graph inductive bias justifies the architectural constraint.

\textbf{Rally-level modelling.} Shot-level classification ignores sequential context. A Transformer over shot embeddings could capture multi-shot strategic arcs.

\begin{thebibliography}{9}
\bibitem{shuttleset}
W.-Y.\ Wang, Y.-C.\ Huang, T.-U.\ Ik, W.-C.\ Peng.
ShuttleSet: A Human-Annotated Stroke-Level Singles Dataset for Badminton Tactical Analysis.
\textit{Proc.\ KDD}, 2023.

\bibitem{bst}
BST: Badminton Stroke-type Transformer for Skeleton-based Action Recognition in Racket Sports.
\textit{arXiv:2502.21085}, 2025.

\bibitem{tempose}
M.\ Ibh, S.\ Gra{\ss}hof, D.~W.\ Hansen.
TemPose: A New Skeleton-Based Transformer Model for Fine-Grained Motion Recognition in Badminton.
\textit{CVPR Workshop on Computer Vision in Sports}, 2023.

\bibitem{shuttlenet}
W.-Y.\ Wang, H.-T.\ Hong, W.-C.\ Peng.
ShuttleNet: Position-Aware Fusion of Rally Progress and Player Styles for Stroke Forecasting in Badminton.
\textit{Proc.\ AAAI}, 2022.

\bibitem{stgcn}
S.\ Yan, Y.\ Xiong, D.\ Lin.
Spatial Temporal Graph Convolutional Networks for Skeleton-Based Action Recognition.
\textit{Proc.\ AAAI}, 2018.
\end{thebibliography}

\end{document}
